\section{A sharper \texorpdfstring{$l$}{l}-truncation bound}
\label{sec:bnd2}

Theorem~\ref{thm:bnd:rem} over-selects $L$ because of two estimates made on the
\emph{singular} solid harmonic: the constant $17$ that counts the terms produced by
$\partial_a\partial_b$, measured loose by a factor $32$ to $34000$, and the
Cauchy--Schwarz $\sum_m|Y_{lm}| \le (2l+1)/\sqrt{4\pi}$. This section replaces both. Two
things change, and each is an exact identity rather than a sharper inequality:

\begin{enumerate}[nosep,leftmargin=1.6em]
\item the two derivatives are moved onto the \emph{weight} by integration by parts. The
triangle weight is an explicit piecewise-linear function whose distributional second
derivative is a bounded measure; the harmonics are then never differentiated at all, and
the factor $\mathrm{hb}(l+2)/\mathrm{hb}(l) \sim (2l+1)(2l+3)/|k\mathbf{R}|^2$ that the old
bound pays for two derivatives on $h_l$ becomes a factor $4/(s_as_b)$;
\item the sum over $m$ is done by \emph{one} Cauchy--Schwarz applied to both factors at
once, $|\sum_m x_my_m| \le \|x\|_2\|y\|_2$, instead of a supremum over $m$ against an
$\ell^1$ norm. Since $\|y\|_2 \le \sqrt{2l+1}\,\|y\|_\infty$, the joint form is never worse
and is typically better by $\sqrt{2l+1}$; and $\sum_m Y_{lm}(\mathbf{u})^2 = (2l+1)/4\pi$
is an equality, so no constant is given away.
\end{enumerate}

\noindent
The scripts are \code{chk0.jl} (the three structural identities), \code{chkA.jl} (the
whole box, end to end), \code{chk0b.jl} and \code{chkB.jl} (the sub-box); the deliverable
is \code{bound2.jl}. Raw output in \code{out\_chk0.txt}, \code{out\_chkA\_main.txt},
\code{out\_chkA\_far.txt}, \code{out\_chk0b.txt}, \code{out\_chkB\_main.txt}.

\subsection{Which truncation is being bounded}
\label{sec:bnd2:which}

Sections~\ref{sec:exp:der} and~\ref{sec:bnd:rem} use two organizations of the same sum.
They are \emph{not} the same term by term. With
$\psi_{lm}(\mathbf{R}) = h_l(k|\mathbf{R}|)Y_{lm}(\hat{\mathbf{R}})$ and
$\varphi_{lm}(\boldsymbol{\delta}) =
j_l(k|\boldsymbol{\delta}|)Y_{lm}(\hat{\boldsymbol{\delta}})$, shell $l$ is
\begin{equation}
\Sigma^{\mathrm{sng}}_l = \sum_m\bigl[(\partial_a\partial_b+\delta_{ab}k^2)\psi_{lm}\bigr]
(\mathbf{R})\!\int_D\! w\,\varphi_{lm},
\qquad
\Sigma^{\mathrm{reg}}_l = \sum_m \psi_{lm}(\mathbf{R})\!\int_D\! w\,
(\partial_a\partial_b+\delta_{ab}k^2)\varphi_{lm},
\label{eq:bnd2:two}
\end{equation}
and $\Sigma^{\mathrm{sng}}_l \neq \Sigma^{\mathrm{reg}}_l$: the first differentiates
$S_l(\mathbf{R},\boldsymbol{\delta})$ in $\mathbf{R}$, the second in
$\boldsymbol{\delta}$, and $S_l$ is not a function of
$\mathbf{R}+\boldsymbol{\delta}$. Both series converge to $T_{ab}$. \emph{The library
truncates the second}, because its geometry tables carry the derivatives
(\eqref{eq:exp:tabs}), while Theorem~\ref{thm:bnd:rem} bounds the first. Measured over
$116$ rows (\code{chkA.jl}), $\max_{ab}|T-T^{(L),\mathrm{reg}}|$ is a median $17.5$ times
$\max_{ab}|T-T^{(L),\mathrm{sng}}|$ at the same $L$, and the $L$ needed for
$10^{-13}\,\mathrm{est}$ is a median one shell larger for the regular organization
($L^{\mathrm{sng}}_{\mathrm{meas}}/L^{\mathrm{reg}}_{\mathrm{meas}}$ has median $0.909$,
range $0.75$--$1$). The old bound is loose enough to cover the difference, so nothing
computed so far is wrong; but a bound meant to \emph{select} $L$ should bound the sum that
is actually truncated. Theorem~\ref{thm:bnd2:reg} does; Theorem~\ref{thm:bnd2:sng} is the
sharpened form of Theorem~\ref{thm:bnd:rem} for the other organization.

\subsection{Three exact identities}
\label{sec:bnd2:lem}

\begin{lemma}[integration by parts onto the triangle weight]
\label{lem:bnd2:ibp}
Extend $w(\boldsymbol{\delta}) = \prod_i(s_i-|\delta_i|)_+$ by zero to $\mathbb{R}^3$. For
every $\phi$ of class $C^2$ on a neighbourhood of $\overline{D}$ and every $a \neq b$,
with $c$ the remaining axis,
\begin{equation}
\int_D w\,\partial_a\partial_b\phi\,\dd\boldsymbol{\delta}
= \int_D \operatorname{sign}(\delta_a)\operatorname{sign}(\delta_b)\,(s_c-|\delta_c|)\,
\phi\,\dd\boldsymbol{\delta},
\label{eq:bnd2:ibpo}
\end{equation}
and for every $a$, writing $Q_a$ for the rectangle in the two axes $b,c$,
\begin{equation}
\int_D w\,\partial_a^2\phi\,\dd\boldsymbol{\delta}
= \int_{Q_a}(s_b-|\delta_b|)(s_c-|\delta_c|)\Bigl[
\phi\big|_{\delta_a = s_a} + \phi\big|_{\delta_a = -s_a}
- 2\,\phi\big|_{\delta_a = 0}\Bigr]\dd\delta_b\,\dd\delta_c .
\label{eq:bnd2:ibpd}
\end{equation}
\end{lemma}

\begin{proof}
$w$ is Lipschitz, vanishes on $\partial D$, and is a product of one-dimensional tent
functions $w_i(t) = (s_i-|t|)_+$ with $w_i' = -\operatorname{sign}(t)$ a.e.\ and
$w_i'' = \delta_{-s_i} - 2\delta_0 + \delta_{s_i}$ in $\mathcal{D}'(\mathbb{R})$. For
$a \neq b$ integrate by parts once in $\delta_a$ and once in $\delta_b$; both boundary
terms vanish because $w = 0$ on $\partial D$, and
$w_a'w_b' = \operatorname{sign}(\delta_a)\operatorname{sign}(\delta_b)$ gives
\eqref{eq:bnd2:ibpo}. For $a = b$ integrate by parts once in $\delta_a$: the boundary term
vanishes and $-\int w_a'\partial_a\phi = \int\operatorname{sign}(\delta_a)\partial_a\phi
= -\int_{-s_a}^{0}\partial_a\phi + \int_{0}^{s_a}\partial_a\phi
= \phi(s_a)+\phi(-s_a)-2\phi(0)$, which is \eqref{eq:bnd2:ibpd}; equivalently it is
$\langle w_a'',\phi\rangle$.
\end{proof}

\begin{lemma}[the two $m$-sums]
\label{lem:bnd2:msum}
For every unit $\mathbf{u}$, $\sum_m Y_{lm}(\mathbf{u})^2 = (2l+1)/(4\pi)$. For every
solution $f$ of the spherical Bessel equation of order $l$ and every
$\mathbf{x} \neq 0$, with $r = |\mathbf{x}|$,
\begin{equation}
\sum_{m}\bigl|\nabla\bigl(f_l(kr)Y_{lm}(\hat{\mathbf{x}})\bigr)\bigr|^2
= \frac{|k|^2}{4\pi}\Bigl[\,l\,|f_{l-1}(kr)|^2 + (l+1)\,|f_{l+1}(kr)|^2\,\Bigr]
=: \frac{|k|^2}{4\pi}\,\mathcal{G}_f(l).
\label{eq:bnd2:grad}
\end{equation}
\end{lemma}

\begin{proof}
The first is the addition theorem at coincident arguments. For the second, the gradient
formula for a scalar multipole field is
\[
\nabla\bigl(f_lY_{lm}\bigr) = k\Bigl[\sqrt{\tfrac{l}{2l+1}}\,f_{l-1}\,
\mathbf{Y}_{l,l-1,m} - \sqrt{\tfrac{l+1}{2l+1}}\,f_{l+1}\,\mathbf{Y}_{l,l+1,m}\Bigr],
\]
$\mathbf{Y}_{Jl'm}$ the vector spherical harmonics, which uses only
$f_l' - (l/z)f_l = -f_{l+1}$ and $f_l' + ((l+1)/z)f_l = f_{l-1}$. The matrix
$M_{l'l''}(\hat{\mathbf{x}}) = \sum_m \mathbf{Y}_{Jl'm}(\hat{\mathbf{x}})\otimes
\overline{\mathbf{Y}_{Jl''m}(\hat{\mathbf{x}})}$ satisfies $M(Q\hat{\mathbf{x}}) =
QM(\hat{\mathbf{x}})Q^{\mathsf{T}}$ for every rotation $Q$, by unitarity of the Wigner
matrices; hence its trace is a rotation-invariant function of $\hat{\mathbf{x}}$, i.e.\ a
constant, and integrating over the sphere with the orthonormality of the
$\mathbf{Y}_{Jl'm}$ gives $\sum_m\mathbf{Y}_{Jl'm}\!\cdot\!\overline{\mathbf{Y}_{Jl''m}}
= \delta_{l'l''}(2J+1)/(4\pi)$. The cross term therefore vanishes and
$l/(2l+1) + (l+1)/(2l+1)$ recombines into \eqref{eq:bnd2:grad}. Passing from complex to
real harmonics is a unitary change of basis inside the shell and leaves every
$\sum_m|\cdot|^2$ unchanged.
\end{proof}

\begin{lemma}[the Hessian, in $\ell^2$ over $m$ --- the replacement of the constant $17$]
\label{lem:bnd2:hess}
With $f$, $l$, $r$ as in Lemma~\ref{lem:bnd2:msum},
\begin{equation}
\Bigl(\sum_{m}\sum_{a,b}\bigl|\partial_a\partial_b\bigl(f_l(kr)Y_{lm}\bigr)\bigr|^2
\Bigr)^{1/2}
\;\le\;
\frac{|k|^{2}}{\sqrt{4\pi}}\left[\sqrt{\frac{l\,\mathcal{G}_f(l-1)}{2l-1}}
+ \sqrt{\frac{(l+1)\,\mathcal{G}_f(l+1)}{2l+3}}\right]
=: \Theta_f(l,r),
\label{eq:bnd2:hess}
\end{equation}
with equality at $l = 0$. Only $f_{l-2}, f_l, f_{l+2}$ enter.
\end{lemma}

\begin{proof}
Apply the gradient formula once: $\partial_b(f_lY_{lm}) = k[a_lf_{l-1}
(\mathbf{Y}_{l,l-1,m})_b - b_lf_{l+1}(\mathbf{Y}_{l,l+1,m})_b]$ with
$a_l = \sqrt{l/(2l+1)}$, $b_l = \sqrt{(l+1)/(2l+1)}$, and use Minkowski's inequality in
$\ell^2(m,a,b)$ to separate the two channels. In the first channel, the assignment
$m \mapsto \bigl((\mathbf{Y}_{l,l-1,m})_b\bigr)_{b=1,2,3}$ is an equivariant map from the
irreducible representation $D^l$ into $D^1\otimes D^{l-1} = D^{l-2}\oplus D^{l-1}\oplus
D^{l}$, in which $D^l$ occurs once; writing $A$ for its matrix in the orthonormal bases
$\{Y_{lm}\}$ and $\{\mathbf{e}_b Y_{l-1,m'}\}$, Schur's lemma gives $AA^\dagger = \lambda
I$ with $\lambda(2l+1) = \operatorname{tr}(AA^\dagger) = \sum_m\int|\mathbf{Y}_{l,l-1,m}|^2
= 2l+1$, so $A$ is a coisometry and $A^\dagger A$ is the orthogonal projector onto that
$D^l$ component. Its partial trace over the vector index is equivariant on $D^{l-1}$, hence
a multiple of the identity, and taking traces the multiple is $(2l+1)/(2l-1)$. Therefore
\[
\sum_{m,b}\bigl|\nabla\bigl(f_{l-1}(\mathbf{Y}_{l,l-1,m})_b\bigr)\bigr|^2
= \frac{2l+1}{2l-1}\sum_{m'}\bigl|\nabla\bigl(f_{l-1}Y_{l-1,m'}\bigr)\bigr|^2
= \frac{2l+1}{2l-1}\cdot\frac{|k|^2}{4\pi}\mathcal{G}_f(l-1),
\]
and $a_l^2(2l+1)/(2l-1) = l/(2l-1)$. The second channel is identical with
$b_l^2(2l+1)/(2l+3) = (l+1)/(2l+3)$. At $l = 0$ there is one channel, so Minkowski is an
equality; direct computation of the Hessian of a radial function confirms
$\Theta_f(0,r)^2 = (|k|^4/4\pi)(|f_0|^2+2|f_2|^2)/3$ on both sides.
\end{proof}

\paragraph{Measured (\code{chk0.jl}, $400$-bit, central differences at step $10^{-30}$).}
\eqref{eq:bnd2:grad} holds as an identity to the last printed digit (ratio
$1.000000$) in $48$ rows, $f \in \{j,h\}$, four points, $l \le 20$.
\eqref{eq:bnd2:hess} is never violated and the ratio true$/$bound is
\[
\begin{array}{lccccc}
|k|r & l=1 & 2 & 5 & 12 & 20\\
0.785\ (h) & 0.958 & 0.984 & 0.996 & 0.9991 & 0.9997\\
0.785\ (j) & 0.708 & 0.945 & 0.993 & 0.9989 & 0.9996\\
9.425\ (h) & 0.719 & 0.712 & 0.710 & 0.849 & 0.949
\end{array}
\]
i.e.\ the replacement of the constant $17$ is loose by at most $1.41$ and typically by less
than $1.01$, against $32$--$34000$ for the constant it replaces.
\eqref{eq:bnd2:ibpo}--\eqref{eq:bnd2:ibpd} were checked against direct differentiation of
$\varphi_{lm}$ under the integral: relative difference $2.5\times10^{-50}$ to
$2.0\times10^{-44}$ (quadrature limited) for $(1/32)^3$ and $(1/32,1/32,1/512)$, all six
$(a,b)$, $l \le 8$.

\subsection{The whole box}
\label{sec:bnd2:whl}

\begin{theorem}[$l$-truncation, derivatives on the regular side]
\label{thm:bnd2:reg}
Let the pair be separated, $\rho < 1$, and let $T^{(L)}_{ab}$ be \eqref{eq:exp:eval}
truncated at $l \le L$. Then for every $a,b$
\begin{equation}
\bigl|T_{ab}-T^{(L)}_{ab}\bigr| \;\le\;
\frac{|k|}{4\pi|f|^{2}V_{\mathrm{t}}}
\sum_{\substack{l>L\\ l\ \mathrm{even}}} (2l+1)\,
\bigl|h^{(1)}_l(k|\mathbf{R}|)\bigr|\;
\frac{|k|^{l}e^{(|k|r_{\mathrm{d}})^{2}/(4l+6)}}{(2l+1)!!}\;V^{ab}_l ,
\label{eq:bnd2:reg}
\end{equation}
where $V^{ab}_l$ is the exact positive geometry moment, for $l = 2p$,
\begin{align}
V^{ab}_l &= \int_D (s_c-|\delta_c|)\,|\boldsymbol{\delta}|^{l}\dd\boldsymbol{\delta}
= \sum_{i+j+q=p}\tfrac{p!}{i!j!q!}\lambda_{2i}(s_a)\lambda_{2j}(s_b)\mu_{2q}(s_c),
&& a\neq b,
\label{eq:bnd2:Vo}\\
V^{aa}_l &= 2\!\!\sum_{i+j+q=p}\!\!\tfrac{p!}{i!j!q!}s_a^{2i}\mu_{2j}(s_b)\mu_{2q}(s_c)
+ 2\!\sum_{j+q=p}\!\tfrac{p!}{j!q!}\mu_{2j}(s_b)\mu_{2q}(s_c) + |k|^2 W_l ,
\label{eq:bnd2:Vd}
\end{align}
with $\lambda_n(s) = \int_{-s}^{s}t^n\dd t = 2s^{n+1}/(n+1)$, $\mu_n$ as in
\eqref{eq:vol:mu} and $W_l$ as in \eqref{eq:bnd:W}. Every term is positive.
\end{theorem}

\begin{proof}
Because $g$ depends on $\mathbf{R}+\boldsymbol{\delta}$ only,
$\partial_{R_a}\partial_{R_b}g = \partial_{\delta_a}\partial_{\delta_b}g$; inserting
\eqref{eq:exp:g}, which converges uniformly with all $\boldsymbol{\delta}$-derivatives on
$\overline{D}$ when $\rho<1$, gives
$T_{ab} = \frac{\ii k}{f^2V_{\mathrm{t}}}\sum_{l,m}(-1)^l\psi_{lm}(\mathbf{R})
\mathcal{M}^{ab}_{lm}$ with $\mathcal{M}^{ab}_{lm} = \int_D w(\partial_a\partial_b +
\delta_{ab}k^2)\varphi_{lm}$, which is $V_{\mathrm{t}}\widetilde{Q}^{ab}_{lm}$ of
\eqref{eq:exp:Q}. Hence
\[
\bigl|T_{ab}-T^{(L)}_{ab}\bigr| \le \frac{|k|}{|f|^2V_{\mathrm{t}}}\sum_{l>L}
\Bigl(\sum_m|\psi_{lm}(\mathbf{R})|^2\Bigr)^{1/2}
\Bigl(\sum_m|\mathcal{M}^{ab}_{lm}|^2\Bigr)^{1/2},
\]
and the first factor is $\sqrt{(2l+1)/4\pi}\,|h_l(k|\mathbf{R}|)|$ exactly by
Lemma~\ref{lem:bnd2:msum}. For the second, apply Lemma~\ref{lem:bnd2:ibp} to
$\phi = \varphi_{lm}$ (entire in $\boldsymbol{\delta}$), then Minkowski's integral
inequality in $\ell^2(m)$, i.e.\
$\|\int F\dd\nu\|_{\ell^2(m)} \le \int\|F\|_{\ell^2(m)}\dd\nu$ for the positive measures
$\nu$ produced by the lemma, together with
$\|\varphi_{lm}(\boldsymbol{\delta})\|_{\ell^2(m)} =
\sqrt{(2l+1)/4\pi}\,|j_l(k|\boldsymbol{\delta}|)|$:
\[
\Bigl(\sum_m|\mathcal{M}^{ab}_{lm}|^2\Bigr)^{1/2} \le
\sqrt{\frac{2l+1}{4\pi}}\;U^{ab}_l,
\]
$U^{ab}_l = \int_D(s_c-|\delta_c|)|j_l(k|\boldsymbol{\delta}|)|$ for $a\neq b$ and, for
$a=b$, $2\int_{Q_a}w_bw_c\bigl[|j_l|(\varrho_+) + |j_l|(\varrho_0)\bigr] +
|k|^2\int_Dw|j_l|$ with $\varrho_+^2 = s_a^2+\delta_b^2+\delta_c^2$ and
$\varrho_0^2 = \delta_b^2+\delta_c^2$. Finally
$|j_l(kx)| \le (|k|x)^le^{(|k|r_{\mathrm{d}})^2/(4l+6)}/(2l+1)!!$ pointwise
(Theorem~\ref{thm:bnd:j}, using $x \le r_{\mathrm{d}}$ on $\overline{D}$ and on both
sections), and integrating $x^l = |\boldsymbol{\delta}|^{2p}$ against those measures
multinomially gives \eqref{eq:bnd2:Vo}--\eqref{eq:bnd2:Vd}. Odd $l$ is omitted because
$\mathcal{M}^{ab}_{lm} = 0$ there (Section~\ref{sec:exp:parity}).
\end{proof}

\begin{remark}[why this is so much smaller]
\label{rem:bnd2:why}
Term by term the ratio of \eqref{eq:bnd2:reg} to \eqref{eq:bnd:rem} is
\[
\frac{|h_l(k\mathbf{R})|\,V^{ab}_l}
{17\,|k|^2\sqrt{2l+5}\;W_l\,\mathrm{hb}(l+2,k\mathbf{R})}
\;\approx\;
\frac{4\,\sigma\,(|\mathbf{R}|/s)^2}{23\,\sqrt{2l+5}\,(2l+1)(2l+3)},
\]
using $V^{ab}_0/W_0 = 4/(s_as_b)$ and $|h_l|/\mathrm{hb}(l+2)\approx
|k\mathbf{R}|^2/((2l+1)(2l+3)\cdot 1.35)$ at small argument, $\sigma$ an $O(1)$ moment
spread. Two derivatives cost $\sim(2l+1)^2/|k\mathbf{R}|^2$ on the singular side and
$\sim 4/(s_as_b)$ on the regular side; the regular side wins whenever $l \gtrsim
|\mathbf{R}|/s$, which is the whole range in which the expansion is used, since
$L \approx 13/\log_{10}(1/(0.7\rho))$ while $|\mathbf{R}|/s = \sqrt3/\rho$.
\end{remark}

\subsection{One affine-weight sub-box}
\label{sec:bnd2:sub}

On a sub-box the weight does not vanish on the boundary, so integration by parts leaves
boundary terms; on the diagonal it leaves one first derivative, which
Lemma~\ref{lem:bnd2:msum} controls exactly.

\begin{theorem}[$l$-truncation on a sub-box]
\label{thm:bnd2:sub}
Let $D_j = \prod_i[-h_i,h_i]$ carry the weight $u(\boldsymbol{\delta}') =
\prod_i u_i(\delta'_i)$, $u_i(t) = \alpha_i+\beta_i t \ge 0$, and let
$\mathbf{V} = \mathbf{R}+\mathbf{c}_j$, $r_j = |\mathbf{h}|$,
$u_i^{\pm} = u_i(\pm h_i) \ge 0$. Then, with all $l$ contributing,
\begin{equation}
\bigl|T^{(j)}_{ab}-T^{(j),L}_{ab}\bigr| \le
\frac{|k|}{4\pi|f|^{2}V_{\mathrm{t}}}\sum_{l>L}(2l+1)\bigl|h^{(1)}_l(k|\mathbf{V}|)\bigr|
\;\Xi^{ab}_l ,
\label{eq:bnd2:sub}
\end{equation}
where, with $\jmath(q) = |k|^{q}e^{(|k|r_j)^2/(4q+6)}/(2q+1)!!$ and $N[\cdot]$ the
$|\boldsymbol{\delta}'|^{q}$-moment of the indicated product measure (odd $q$ by
Cauchy--Schwarz, $N_q \le \sqrt{N_{q-1}N_{q+1}}$),
\begin{align}
\Xi^{ab}_l &= \jmath(l)\,N_l\bigl[\nu_a\otimes\nu_b\otimes u_c\bigr], && a \ne b,\\
\Xi^{aa}_l &= \sqrt{\tfrac{l}{2l+1}}\,\jmath(l-1)\,|k|\,N_{l-1}[\pi_a\otimes u_b\otimes u_c]
+ \sqrt{\tfrac{l+1}{2l+1}}\,\jmath(l+1)\,|k|\,N_{l+1}[\pi_a\otimes u_b\otimes u_c]
\notag\\
&\qquad + \jmath(l)\,N_l[\beta_a\otimes u_b\otimes u_c]
+ |k|^2\jmath(l)\,N_l[u_a\otimes u_b\otimes u_c],
\end{align}
and the one-dimensional positive measures are
$\nu_i = u_i^{+}\delta_{h_i}+u_i^{-}\delta_{-h_i}+|\beta_i|\dd t$,
$\pi_i = u_i^{+}\delta_{h_i}+u_i^{-}\delta_{-h_i}$,
$\beta_i = |\beta_i|(\delta_{h_i}+\delta_{-h_i})$.
\end{theorem}

\begin{proof}
The one-dimensional identities are
$\int_{-h}^{h}u\,X' = u^{+}X(h) - u^{-}X(-h) - \beta\int_{-h}^{h}X$ and
$\int_{-h}^{h}u\,X'' = u^{+}X'(h) - u^{-}X'(-h) - \beta\bigl(X(h)-X(-h)\bigr)$. For
$a \ne b$ apply the first in axis $a$ and in axis $b$ and keep $u_c$ in axis $c$; taking
absolute values term by term replaces each one-dimensional functional by the positive
measure $\nu_i$, and Minkowski in $\ell^2(m)$ with
$\|\varphi_{lm}\|_{\ell^2(m)} = \sqrt{(2l+1)/4\pi}|j_l|$ gives the first line. For $a = b$
apply the second identity in axis $a$: the two surviving derivative evaluations are
bounded by $\|\nabla\varphi_{lm}\|$, and Lemma~\ref{lem:bnd2:msum} with
$|j_{l\mp1}| \le \jmath(l\mp1)|\boldsymbol{\delta}'|^{l\mp1}$ and
$\sqrt{A+B}\le\sqrt A+\sqrt B$ gives the first two terms; the $\beta_a$ term and the
$\delta_{ab}k^2$ term give the last two. Every measure is positive, so every moment is an
exact positive multinomial sum in the one-dimensional moments
$\int|t|^{2i}\dd\nu$, and $\rho_j<1$ makes the interchange of sum and integral legitimate.
For the octant split $u_i^{+} = 0$ and $u_i^{-} = s_i$, so only the inner face carries a
derivative.
\end{proof}

\subsection{The singular organization, with the constant \texorpdfstring{$17$}{17} removed}
\label{sec:bnd2:sng}

\begin{theorem}[sharpened form of Theorem~\ref{thm:bnd:rem}]
\label{thm:bnd2:sng}
With $T^{(L),\mathrm{sng}}$ the truncation of \eqref{eq:exp:sing} and
$\Theta_h(l,|\mathbf{R}|)$ of Lemma~\ref{lem:bnd2:hess} built on $f = h^{(1)}$,
\begin{equation}
\bigl|T_{ab}-T^{(L),\mathrm{sng}}_{ab}\bigr| \le
\frac{|k|}{|f|^{2}V_{\mathrm{t}}}\sum_{\substack{l>L\\ l\ \mathrm{even}}}
\left[\Theta_h(l,|\mathbf{R}|)+\delta_{ab}|k|^2\sqrt{\tfrac{2l+1}{4\pi}}|h_l|\right]
\sqrt{\tfrac{2l+1}{4\pi}}\;
\frac{|k|^{l}e^{(|k|r_{\mathrm{d}})^{2}/(4l+6)}}{(2l+1)!!}\,W_l .
\label{eq:bnd2:sng}
\end{equation}
\end{theorem}

\begin{proof}
As in Theorem~\ref{thm:bnd2:reg} but with the roles exchanged: Cauchy--Schwarz over $m$
between $\bigl((\partial_a\partial_b+\delta_{ab}k^2)\psi_{lm}\bigr)_m$ and
$(\mathcal{I}_{lm})_m$; the first $\ell^2$ norm is at most
$\Theta_h + \delta_{ab}|k|^2(\sum_m|\psi_{lm}|^2)^{1/2}$ by Minkowski and
Lemma~\ref{lem:bnd2:hess}, the second is at most
$\sqrt{(2l+1)/4\pi}\int_Dw|j_l|$ by Minkowski and Lemma~\ref{lem:bnd2:msum}.
\end{proof}

\noindent
Against \eqref{eq:bnd:rem} the term-by-term ratio is
$\Theta_h\sqrt{(2l+1)/4\pi}\big/\bigl(17|k|^2\mathrm{hb}(l+2)\sqrt{(2l+5)/4\pi}\,
(2l+1)/\sqrt{4\pi}\bigr) \approx 1/(34\sqrt{2l}\cdot 1.35)$, i.e.\ the improvement grows
like $\sqrt{l}$: the old bound loses a factor $\sqrt{2l+5}$ that the exact $m$-sum does
not.

\subsection{Two routes that do not help, and one that helps a little}
\label{sec:bnd2:routes}

\paragraph{A Cauchy estimate in one complex variable.} Along
$\mathbf{R}+t\boldsymbol{\delta}$ the kernel is analytic for $|t| < |\mathbf{R}|/
|\boldsymbol{\delta}|$ (Proposition~\ref{prop:bnd:rho}) and a Cauchy estimate bounds the
tail of its \emph{Taylor} series in $t$. That is the degree-graded truncation of family
(a), not the $l$-truncation: shell $l$ contains the degrees $l, l+2, l+4,\dots$ and degree
$q$ contains the shells $q, q-2,\dots$, so neither tail majorises the other. The estimate
therefore cannot be substituted here, and where it does apply it is $2.0$--$7.7$ times
pessimistic in the selected degree.

\paragraph{Exact Bessel magnitudes instead of majorants.} Both $|h_l|$ and $|j_l|$ are
computable to full precision at negligible cost, and using them in place of the majorants
of Theorems~\ref{thm:bnd:j} and~\ref{thm:bnd:h} keeps the bound rigorous up to the
rounding of an elementary function. Measured (\code{chkA.jl}, $29$ cases, tails at
$L = 32$): replacing $|h_l|$ by $\mathrm{hb}(l,k\mathbf{R})$ inflates the bound by $1.48$
at $|k\mathbf{R}| = 0.39$, $2.17$ at $0.79$, $4.6$ at $1.6$, $13.6$ at $2.7$, $20$ at $3.1$,
$297$ at $6.3$, $2.7\times10^{4}$ at $12.6$ and $2.3\times10^{6}$ at $21.8$ --- the regime
$|k\mathbf{R}| \gtrsim l$, where the triangle inequality on the finite Hankel sum gives
away the oscillation --- but costs only a median $8\%$ (at most $38\%$) in the selected
$L$, because $L$ is set by $\rho$ and a constant factor moves it by a few shells.
Replacing the exponential majorant of $j_l$ by the termwise-absolute power series
$\sum_n |k|^{l+2n}V_{l+2n}/(2^nn!(2l+2n+1)!!)$ gains between $0.002\%$ and $2.2\%$ and
never changes the selected $L$: the $j_l$ majorant is not where the looseness is.

\subsection{Measured}
\label{sec:bnd2:meas}

Truth: $\max_{ab}|T_{ab}-T^{(L)}_{ab}|$ for the organization named, with $T_{ab}$ from an
independent \code{BigFloat} volume Gauss--Legendre rule on the eight octants of $D$
(order $22$; agreement with the shell sums at $L = 32$ is $10^{-18}$ to $10^{-57}$
relative) and the shell data from a $24$-point positive-octant rule with exact parity
factors. $116$ rows: four cell shapes, offsets from $2$ to $64$ cells,
$f \in \{1, 1+0.1\ii, 1+\ii\}$, $L \in \{8,16,24,32\}$; $320$-bit, and $512$-bit with a
$10^{-38}$ difference step for the far offsets, where the central-difference evaluation of
$(\partial_a\partial_b+\delta_{ab}k^2)\psi_{lm}$ has a floor. No violation.

\begin{center}\footnotesize
\begin{tabular}{@{}lrrrr@{}}
\toprule
bound / truth & min & median & max \\
\midrule
\eqref{eq:bnd2:reg}, regular organization & 8.79 & 49.1 & $2.93\times10^{3}$\\
\eqref{eq:bnd2:sng}, singular organization & 1.49 & 140 & 792\\
\eqref{eq:bnd:rem}, the first-round bound & $1.35\times10^{3}$ & $1.53\times10^{5}$ &
$1.73\times10^{11}$\\
\midrule
gain \eqref{eq:bnd:rem}/\eqref{eq:bnd2:reg} & 31.3 & $3.23\times10^{3}$ &
$4.61\times10^{9}$\\
gain \eqref{eq:bnd:rem}/\eqref{eq:bnd2:sng} & 162 & 661 & $7.21\times10^{8}$\\
\bottomrule
\end{tabular}
\end{center}

\noindent
Cost: $V^{ab}_l$ depends on the cell shape alone once the $|k|^2W_l$ piece is split off, so it
is tabulated once per shape ($31$~ms to $l = 140$) and the per-offset work is the $h_l$
recurrence and one length-$70$ dot product, measured at $3.8$, $3.9$ and $9.7\,\mu$s at
$L = 8$, $16$, $32$ on a machine that was not quiet.

\noindent
The largest ratios of \eqref{eq:bnd2:reg} are not looseness of the bound but dips of the
truth: the true remainder oscillates in $L$ (the $l$-sum alternates with the phase of
$h_l$), and no monotone majorant can follow a dip. At the $L$ that matters --- the one the
bound selects --- the ratio is $9$--$100$.

\begin{center}\footnotesize
\begin{tabular}{@{}llrrrrrr@{}}
\toprule
shape & $\mathbf{n}$ & $\rho$ & $|k\mathbf{R}|$ & $L_{\mathrm{meas}}$ &
$L$ from \eqref{eq:bnd2:reg} & $L$ from \eqref{eq:bnd:rem} & cost ratio\\
\midrule
$(1/32)^3$ & $(2,0,0)$  & 0.866 & 0.39  & $>32$ & 126 & $>136$ & --\\
$(1/32)^3$ & $(3,0,0)$  & 0.577 & 0.59  & 32 & 40 & 50 & 1.56\\
$(1/32)^3$ & $(4,0,0)$  & 0.433 & 0.79  & 24 & 28 & 34 & 1.47\\
$(1/32)^3$ & $(3,1,0)$  & 0.548 & 0.62  & 32 & 36 & 46 & 1.63\\
$(1/32)^3$ & $(8,0,0)$  & 0.217 & 1.57  & 16 & 16 & 20 & 1.56\\
$(1/32)^3$ & $(8,8,8)$  & 0.125 & 2.72  & 12 & 12 & 14 & 1.36\\
$(1/32)^3$ & $(2,2,0)$  & 0.612 & 0.56  & $>32$ & 44 & 56 & 1.62\\
$(1/8)^3$  & $(4,0,0)$  & 0.433 & 3.14  & 22 & 26 & 36 & 1.92\\
$(1/8)^3$  & $(3,1,0)$  & 0.548 & 2.50  & 28 & 36 & 48 & 1.78\\
$(1/4)^3$  & $(4,0,0)$  & 0.433 & 6.28  & 22 & 26 & 38 & 2.14\\
$(1/4)^3$  & $(8,0,0)$  & 0.217 & 12.57 & 14 & 16 & 24 & 2.25\\
$(1/4)^3$  & $(8,8,8)$  & 0.125 & 21.77 & 14 & 16 & 20 & 1.56\\
sl         & $(4,0,0)$  & 0.354 & 0.79  & 22 & 26 & 30 & 1.33\\
sl         & $(0,0,64)$ & 0.354 & 0.79  & 24 & 26 & 30 & 1.33\\
sl         & $(0,0,32)$ & 0.708 & 0.39  & $>32$ & 70 & 90 & 1.65\\
\bottomrule
\end{tabular}
\end{center}

\noindent
$L_{\mathrm{meas}}$ is the smallest even $L$ with
$\max_{ab}|T-T^{(L),\mathrm{reg}}| \le 10^{-13}\mathrm{est}(\mathbf{R})$; ``cost ratio'' is
$(L_{\mathrm{old}}/L_{\mathrm{new}})^2$, the saving in the $O(L^2)$ per-offset work. Over
$25$ cases with all three defined, $L_{\mathrm{old}}/L_{\mathrm{new}}$ is $1.00$--$1.50$
with median $1.21$ (cost $1.00$--$2.25$, median $1.47$) and the new bound over-selects by
$L_{\mathrm{new}}/L_{\mathrm{meas}} = 1.00$--$1.29$, median $1.14$, against
$1.17$--$1.73$, median $1.40$, for the bound it replaces. At $(2,0,0)$ the old bound gives
no certificate below $L = 136$ while \eqref{eq:bnd2:reg} certifies $L = 126$.

\subsection{What is still loose}
\label{sec:bnd2:loose}

The remaining factor of $9$--$100$ is, in order of size: (i) the Cauchy--Schwarz over $m$,
which is an equality only when $\mathcal{M}_{lm}$ is proportional to $Y_{lm}(\hat{\mathbf
R})$ --- for the whole box only one parity class per entry is non-zero, so a factor between
$1$ and $\sqrt{\#\{m\ \mathrm{surviving}\}}$ is given away and is not recovered by any
argument that does not know the geometry table; (ii) Minkowski under the
$\boldsymbol{\delta}$ integral, which discards the cancellation of $Y_{lm}$ over the box
--- the same loss the $W_l$ moment already minimises for the radial part; (iii) the
absolute values across $l$, worth the measured $\sum_l|t_l|/|T| \le 8.4$ at $\lambda/32$
and $\le 67$ at $\lambda/4$; (iv) $1.00$--$1.41$ for Lemma~\ref{lem:bnd2:hess} and
$1.0006$--$1.5$ for the $j_l$ majorant. Nothing here bounds a rounding error, and
\eqref{eq:bnd2:reg} uses $|h_l|$ from the upward recurrence, which is a measurement
($3\times10^{-15}$ for $l \le 40$) and not a theorem; the variant with
$\mathrm{hb}(l,k\mathbf{R})$ is fully certified and costs a median $8\%$ in $L$.
